%% ==========================================================================
%%  Appendix C  The one-good extension when the bundles contain chores
%%
%%  Full check of Lemma~\ref{lem:goods-extension}, moved out of Section 4 on
%%  2026-10-08 (Section 4 now cites the algorithm of Barman et al. directly).
%%  Text of arXiv Appendix B (arxiv/sections/appendix_goods_extension.tex) in
%%  the house style, plus one added sentence in Case 2: s_{tau+1} differs from
%%  s_tau, so s_0 has a predecessor on Q_1. Lemma 10 of Barman et al. omits the
%%  same case.
%% ==========================================================================
\section{The One-Good Extension when the Bundles Contain Chores}
\label{app:goods-extension}

This appendix restates the one-good extension argument of Barman et
al.~\cite{BKNS22} in a self-contained way and checks that each step uses only
the two facts named in the proof of Lemma~\ref{lem:goods-extension}: the
valuations are integer valued, and the good being inserted has marginal value in
$\{0,1\}$ for every agent on every bundle. In particular, no step uses the sign
of the items already allocated.

\begin{proof}[Proof of Lemma~\ref{lem:goods-extension}]
For an allocation $Y=(Y_1,\ldots,Y_n)$, let $G_Y$ denote its weighted envy
graph, with edge weights
\[
  w_Y(i,j):=v_i(Y_j)-v_i(Y_i).
\]
We use two standard characterizations of envy-freeability~\cite{HS19,BKNS22}
(Theorems~\ref{thm:hs-characterisation} and~\ref{thm:hs-minsubsidy}). First,
$Y$ is envy-freeable if and only if it maximizes
\[
  \sum_{i\in N}v_i(Y_{\sigma(i)})
\]
over all permutations $\sigma$ of its bundles. Equivalently, $G_Y$ contains no
positive-weight directed cycle. Second, if $Y$ is envy-freeable, then its
pointwise-minimal supporting subsidy for agent $i$ is the maximum weight of a
directed path starting at $i$ in $G_Y$.

We may assume without loss of generality that $p$ is the pointwise-minimal
subsidy vector supporting $A$. Indeed, starting from the subsidy vector in the
statement, replace it by the pointwise-minimal one. This can only decrease the
subsidies componentwise. Since the valuations are integer-valued, all edge
weights in $G_A$, and therefore all pointwise-minimal subsidies, are integers.
Consequently, the resulting vector still belongs to $\{0,1\}^n$.

We first prove two properties of adding $g$ that will be used throughout the
proof.

\medskip
\noindent
\textbf{Claim 1.}
Let $Y$ be an envy-freeable allocation of $M'$, and suppose that for some agent
$x$,
\[
  v_x(Y_x\cup\{g\})-v_x(Y_x)=1.
\]
Then
\[
  Z=(Y_1,\ldots,Y_x\cup\{g\},\ldots,Y_n)
\]
is envy-freeable.

\smallskip
\noindent
\emph{Proof of Claim 1.}
Write
\[
  \operatorname{SW}(Y):=\sum_{i\in N}v_i(Y_i).
\]
Since the marginal value of $g$ for $x$ on $Y_x$ is one,
\[
  \operatorname{SW}(Z)=\operatorname{SW}(Y)+1.
\]
Consider any permutation $\sigma$ of the bundles of $Z$. Let $a$ be the agent who
receives the bundle containing $g$ under this permutation. Removing $g$ from that
bundle leaves a permutation of the bundles of $Y$. Since $Y$ is envy-freeable,
the welfare of this reassignment is at most $\operatorname{SW}(Y)$. Moreover,
\[
  v_a(S\cup\{g\})-v_a(S)\leq 1
\]
for every $S\subseteq M'$. Hence
\[
  \operatorname{SW}(Z_\sigma)
  \leq \operatorname{SW}(Y)+1
  =\operatorname{SW}(Z).
\]
Thus $Z$ maximizes welfare over all permutations of its bundles and is therefore
envy-freeable.
\hfill$\square$

\medskip
\noindent
\textbf{Claim 2.}
Let $Y$ be an allocation and let
\[
  Z=(Y_1,\ldots,Y_x\cup\{g\},\ldots,Y_n).
\]
Then
\[
  w_Z(i,j)=w_Y(i,j)
  \qquad
  \text{for all }i,j\neq x,
\]
while
\[
  w_Z(x,j)\leq w_Y(x,j)
  \qquad\text{and}\qquad
  w_Z(i,x)\leq w_Y(i,x)+1.
\]
Consequently, for every simple directed path $P$,
\begin{equation}\label{eq:ge-S1}
  w_Z(P)\leq w_Y(P)+1.
\end{equation}

\smallskip
\noindent
\emph{Proof of Claim 2.}
Only the bundle of $x$ changes. Hence, edges not incident to $x$ are unchanged.
For every $j\neq x$,
\[
\begin{aligned}
  w_Z(x,j)
  &=v_x(Y_j)-v_x(Y_x\cup\{g\})\\
  &=w_Y(x,j)
    -\bigl(v_x(Y_x\cup\{g\})-v_x(Y_x)\bigr)
  \leq w_Y(x,j),
\end{aligned}
\]
because the marginal value of $g$ is nonnegative. Similarly, for $i\neq x$,
\[
\begin{aligned}
  w_Z(i,x)
  &=v_i(Y_x\cup\{g\})-v_i(Y_i)\\
  &=w_Y(i,x)
    +\bigl(v_i(Y_x\cup\{g\})-v_i(Y_x)\bigr)
  \leq w_Y(i,x)+1.
\end{aligned}
\]
A simple path enters $x$ at most once. Its weight can therefore increase by at
most one, proving \eqref{eq:ge-S1}.
\hfill$\square$

\medskip
We now follow the extendable/non-extendable distinction of Barman et
al.~\cite{BKNS22}. For a subsidy vector $q$, recall that
\[
  \Mset(q):=\{i\in N:q_i\geq q_j\text{ for every }j\in N\}.
\]
Call $(A,p)$ \emph{extendable with $g$} if there exists a permutation $\sigma$
such that, writing
\[
  B_i:=A_{\sigma(i)}
  \qquad\text{and}\qquad
  q_i:=p_{\sigma(i)},
\]
the pair $(B,q)$ is envy-free and there is some $\kappa\in \Mset(q)$ satisfying
\begin{equation}\label{eq:ge-S2}
  v_\kappa(B_\kappa\cup\{g\})-v_\kappa(B_\kappa)=1.
\end{equation}

We will also use the following elementary permutation property. If $(A,p)$ is
envy-free and $B=A_\sigma$ is envy-freeable, then $(B,p_\sigma)$ is envy-free.
To see this, the envy-freeness of $(A,p)$ gives
\[
  v_i(A_i)+p_i
  \geq
  v_i(A_{\sigma(i)})+p_{\sigma(i)}
  \qquad\text{for every }i.
\]
Summing these inequalities and using the fact that both $A$ and $B$ are
envy-freeable shows that equality must hold for every $i$. The envy-freeness
inequalities for $(B,p_\sigma)$ then follow directly. See Lemma~3 in Barman et
al.~\cite{BKNS22} for a detailed argument.

\medskip
\noindent
\textbf{Case 1: $(A,p)$ is extendable with $g$.}

Let $\sigma$ and $\kappa$ satisfy the above conditions, and define
\[
  C=(B_1,\ldots,B_\kappa\cup\{g\},\ldots,B_n).
\]
By Claim~1, $C$ is envy-freeable.

It remains to show that $C$ can be supported by subsidies in $\{0,1\}^n$. Since
$q_\kappa$ is maximal and $q\in\{0,1\}^n$, there are two cases.

If $q_\kappa=0$, then $q_i=0$ for all $i\in N$. Thus, every directed path in
$G_B$ has weight at most zero. By Claim~2, every directed path in $G_C$ has
weight at most one. Hence, the pointwise-minimal supporting subsidies for $C$
are at most one. They are nonnegative integers, so they belong to $\{0,1\}^n$.

Suppose next that $q_\kappa=1$. Define
\[
  \widehat q_\kappa:=0,
  \qquad
  \widehat q_i:=q_i
  \quad (i\neq\kappa).
\]
We claim that $(C,\widehat q)$ is envy-free. For $i,j\neq\kappa$, nothing changes
from $(B,q)$. For agent $\kappa$, using \eqref{eq:ge-S2},
\[
  v_\kappa(C_\kappa)+\widehat q_\kappa
  =
  v_\kappa(B_\kappa)+1
  =
  v_\kappa(B_\kappa)+q_\kappa
  \geq
  v_\kappa(B_j)+q_j.
\]
Finally, for every $i\neq\kappa$,
\[
\begin{aligned}
  v_i(C_i)+\widehat q_i
  &=v_i(B_i)+q_i\\
  &\geq v_i(B_\kappa)+q_\kappa\\
  &=v_i(B_\kappa)+1\\
  &\geq v_i(B_\kappa\cup\{g\})
   =v_i(C_\kappa)+\widehat q_\kappa,
\end{aligned}
\]
where the last inequality follows because the marginal value of $g$ is at most
one. Thus $\widehat q\in\{0,1\}^n$ supports $C$.

Hence, the desired extension exists in the extendable case.

\medskip
\noindent
\textbf{Case 2: $(A,p)$ is not extendable with $g$.}

We use the \textsc{FindSink} procedure of Barman et al.~\cite{BKNS22}, and
verify that its proof remains valid in the present setting.

For $s\in N$, write
\[
  X^s:=(A_1,\ldots,A_s\cup\{g\},\ldots,A_n).
\]
Start with an arbitrary $s_0\in \Mset(p)$. Whenever $X^{s_\tau}$ is
envy-freeable, let $\phi^\tau$ denote its pointwise-minimal supporting subsidy
vector. If every component of $\phi^\tau$ is strictly smaller than $2$,
terminate. Otherwise, choose an agent $s_{\tau+1}$ with
\[
  \phi^\tau_{s_{\tau+1}}\geq2
\]
and continue with $X^{s_{\tau+1}}$.

We first establish the following:
\begin{equation}\label{eq:ge-S3}
  s_\tau\in \Mset(p)
  \quad\text{and}\quad
  X^{s_\tau}\text{ is envy-freeable}
\end{equation}
for every agent considered by the procedure.

For $\tau=0$, the first part follows from the choice of $s_0$. Suppose, towards
a contradiction, that $X^{s_0}$ is not envy-freeable. By the welfare
characterization, there exists a permutation of its bundles whose welfare is
strictly larger. Since all valuations are integer-valued, the improvement is at
least one.

Let $B$ be the corresponding permutation of the bundles of $A$, obtained by
deleting $g$ from the permuted bundle containing it, and let $k$ be the agent who
receives the bundle $A_{s_0}$ in $B$. Since the marginal value of $g$ is
nonnegative,
\[
  \operatorname{SW}(X^{s_0})\geq\operatorname{SW}(A).
\]
Therefore,
\begin{equation}\label{eq:ge-S4}
  \operatorname{SW}(B)
  +
  \bigl(v_k(B_k\cup\{g\})-v_k(B_k)\bigr)
  \geq
  \operatorname{SW}(A)+1.
\end{equation}
On the other hand, $A$ is envy-freeable, and hence
\[
  \operatorname{SW}(B)\leq\operatorname{SW}(A).
\]
The marginal term in \eqref{eq:ge-S4} is at most one. Consequently, both
inequalities must be tight:
\[
  \operatorname{SW}(B)=\operatorname{SW}(A)
\]
and
\[
  v_k(B_k\cup\{g\})-v_k(B_k)=1.
\]
Thus $B$ is envy-freeable. By the permutation property above, if $q$ is obtained
by permuting $p$ in the same way, then $(B,q)$ is envy-free. Moreover,
$B_k=A_{s_0}$ and $s_0\in \Mset(p)$, so $k\in \Mset(q)$. Hence $(A,p)$ is
extendable with $g$, a contradiction. Therefore $X^{s_0}$ is envy-freeable.

Now suppose \eqref{eq:ge-S3} holds for $s_\tau$ and the procedure selects
$s_{\tau+1}$ because $\phi^\tau_{s_{\tau+1}}\geq2$. We claim that
$s_{\tau+1}\in \Mset(p)$. If not, then the maximal subsidy in $p$ is one and
$p_{s_{\tau+1}}=0$. Since $p$ is pointwise-minimal, every path starting from
$s_{\tau+1}$ in $G_A$ has weight at most zero. By Claim~2, every path starting
from $s_{\tau+1}$ in $G_{X^{s_\tau}}$ has weight at most one. Hence
$\phi^\tau_{s_{\tau+1}}\leq1$, a contradiction. Thus $s_{\tau+1}\in \Mset(p)$.
The same argument used for $s_0$ then shows that $X^{s_{\tau+1}}$ is
envy-freeable. This proves \eqref{eq:ge-S3}.

Note also that $s_{\tau+1}\neq s_\tau$. Indeed, a simple path starting at
$s_\tau$ never enters $s_\tau$, so by Claim~2 its weight in $G_{X^{s_\tau}}$ is at
most its weight in $G_A$, which is at most $p_{s_\tau}\leq1$. Hence
$\phi^\tau_{s_\tau}\leq1$, and $s_\tau$ is never selected immediately after
itself.

It remains to show that the procedure terminates. Suppose otherwise that some
agent is selected twice. Consider two consecutive occurrences of a repeated agent
and relabel the corresponding segment as
\[
  s_0,s_1,\ldots,s_T,
  \qquad
  s_T=s_0,
\]
with $s_0,\ldots,s_{T-1}$ distinct. By the preceding paragraph, $T\geq2$.

For every $\tau=1,\ldots,T$, agent $s_\tau$ is selected because its
pointwise-minimal subsidy for $X^{s_{\tau-1}}$ is at least two. Hence, there
exists a simple directed path $Q_\tau$ starting at $s_\tau$ in
$G_{X^{s_{\tau-1}}}$ such that
\begin{equation}\label{eq:ge-S6}
  w_{X^{s_{\tau-1}}}(Q_\tau)\geq2.
\end{equation}
The path $Q_\tau$ must contain $s_{\tau-1}$. Otherwise, none of its edge weights
differs from the corresponding edge weights in $G_A$, whereas every path starting
at $s_\tau$ in $G_A$ has weight at most $1$, contradicting \eqref{eq:ge-S6}.

Let $P_\tau$ be the initial segment of $Q_\tau$ from $s_\tau$ to $s_{\tau-1}$. By
Claim~2,
\begin{equation}\label{eq:ge-S7}
  w_A(Q_\tau)
  \geq
  w_{X^{s_{\tau-1}}}(Q_\tau)-1
  \geq1.
\end{equation}
The part of $Q_\tau$ following $P_\tau$ is a path starting at $s_{\tau-1}$ in
$G_A$, and therefore has weight at most $p_{s_{\tau-1}}\leq1$. It follows from
\eqref{eq:ge-S7} that
\begin{equation}\label{eq:ge-S8}
  w_A(P_\tau)\geq0.
\end{equation}

Now concatenate the paths $P_T,P_{T-1},\ldots,P_1$. Since $s_T=s_0$, this
concatenation starts and ends at $s_0$. If some vertex is repeated in the
concatenation, split the resulting closed sequence of edges at repeated vertices.
In this way, it can be decomposed into directed cycles of $G_A$. By
\eqref{eq:ge-S8}, the total weight of the concatenated sequence is nonnegative.
On the other hand, since $A$ is envy-freeable, $G_A$ contains no positive-weight
directed cycle. Therefore, every cycle arising in the above decomposition has
weight at most zero, while their total weight is nonnegative. It follows that the
total weight is zero, and consequently, every cycle in the decomposition has
weight zero.

We now derive a contradiction to non-extendability. Since $s_1$ was selected
after considering $X^{s_0}$, there is a path $Q_1$ satisfying \eqref{eq:ge-S6}.
Its weight in $G_A$ is at most $p_{s_1}\leq1$, while Claim~2 says that its weight
can increase by at most one after $g$ is added to $A_{s_0}$. Hence
\[
  w_{X^{s_0}}(Q_1)=2,
  \qquad
  w_A(Q_1)=1,
\]
and the weight of $Q_1$ increases by exactly one. Only an edge entering $s_0$ can
increase when $g$ is added to $A_{s_0}$. Edges leaving $s_0$ can only decrease.
Since $s_1\neq s_0$ and $Q_1$ starts at $s_1$ and contains $s_0$, the agent $s_0$
has a predecessor $k$ on $Q_1$, and
\begin{equation}\label{eq:ge-S9}
  v_k(A_{s_0}\cup\{g\})-v_k(A_{s_0})=1.
\end{equation}
Since the edge $(k,s_0)$ belongs to $P_1$, it belongs to one of the zero-weight
directed cycles obtained above. Let $C$ be such a cycle. Cyclically permuting the
bundles of $A$ along $C$ preserves social welfare, because the change in welfare
is exactly $w_A(C)=0$. Let $\sigma$ denote the resulting permutation and let
$B=A_\sigma$. Hence
\[
  \operatorname{SW}(B)=\operatorname{SW}(A),
\]
so $B$ is envy-freeable. By the permutation property, $(B,p_\sigma)$ is
envy-free.

Furthermore, the cycle contains $(k,s_0)$, so agent $k$ receives bundle $A_{s_0}$
in $B$. Thus \eqref{eq:ge-S9} gives
\[
  v_k(B_k\cup\{g\})-v_k(B_k)=1.
\]
Since $s_0\in \Mset(p)$, agent $k$ belongs to $\Mset(p_\sigma)$. Consequently,
$\sigma$ and $k$ witness that $(A,p)$ is extendable with $g$, contradicting the
assumption of Case~2.

Therefore, no agent is selected twice, and the procedure terminates after at most
$n$ selections. At termination, the current allocation $X^s$ is envy-freeable by
\eqref{eq:ge-S3}, and every component of its pointwise-minimal subsidy vector is
strictly smaller than $2$. These subsidies are nonnegative integers, and hence
they belong to $\{0,1\}^n$. This proves the existence of the desired allocation
$A'$ and subsidy vector $p'$ in the non-extendable case.

Finally, all steps are polynomial-time. Pointwise-minimal subsidies can be
computed from maximum-weight paths in an envy-freeable allocation. Extendability
can be tested by the \textsc{Extend} subroutine of Barman et al.~\cite{BKNS22}:
for each candidate recipient and most-subsidized bundle, it solves a
maximum-weight bipartite matching problem. Its correctness uses only the welfare
characterization of envy-freeability and the values of the existing bundles, and
is therefore unchanged here. In the non-extendable case, the procedure above
considers at most $n$ agents. Thus, the entire one-good extension can be carried
out in polynomial time using value-oracle access.
\end{proof}
